\documentclass[11pt,a4paper]{article}
\usepackage[utf8]{inputenc}
\usepackage[T1]{fontenc}
\usepackage{amsfonts}
\usepackage[french]{babel}
\frenchbsetup{StandardLists=true}
\usepackage{enumitem}
\usepackage{amssymb}
\usepackage{amsmath}
\usepackage[left=2cm,right=2cm,top=2cm,bottom=2cm]{geometry}
\usepackage{multirow}
\usepackage[dvipsnames]{xcolor}
\usepackage{fouriernc}
\usepackage{graphicx}
\usepackage{setspace}
\singlespacing
\usepackage{tcolorbox}
\usepackage{multicol}
\usepackage{caption}
\usepackage{fancyhdr}
\pagestyle{fancy}
\renewcommand\headrulewidth{1pt}
\fancyhead[L]{Lycée Denis Diderot}
\fancyhead[R]{Classe de 3\textsuperscript{ème}}
\setcounter{page}{1}\fancyfoot[C]{page \thepage}
\fancyfoot[R]{Année 2025-2026}
\fancyfoot[L]{Alexis RUIZ}
\usepackage{multirow,tabularx,booktabs}
\usepackage{colortbl}
\newcommand{\cadreblanc}[1]{%
\medskip\framebox[\linewidth][c]{%
\rule{0mm}{#1mm}}\par
\medskip}
\usepackage{awesomebox}
\setlength{\aweboxvskip}{0mm}
\usepackage{pifont}
\setlength{\parindent}{0pt}
\setlength{\headheight}{14pt}
\usepackage{tikz}
\usetikzlibrary{babel}
\usepackage[european]{circuitikz}
\renewcommand{\arraystretch}{1.8}
\frenchbsetup{StandardLists=true}

% Boîte orange "Appeler le professeur"
\newcommand{\appelprof}{%
\begin{tcolorbox}[colback=orange!10!white, colframe=orange!80!black]
\textbf{Appeler le professeur.}
\end{tcolorbox}}

\begin{document}

%------------------------------------------------------------
% TITRE
%------------------------------------------------------------
\begin{center}
  {\LARGE\textbf{TP -- Schématisation et montages électriques}}
\end{center}

\hrule
\vspace{3mm}
\noindent Nom : \dotfill \quad Prénom : \dotfill
\vspace{3mm}
\hrule
\vspace{5mm}

\begin{tcolorbox}[colback=blue!5!white, colframe=blue!60!black, title=\textbf{Objectifs}]
\begin{itemize}
  \item Connaître et utiliser les symboles normalisés des composants électriques.
  \item Lire et réaliser un schéma électrique.
  \item Distinguer un montage en \textbf{série} d'un montage en \textbf{dérivation}.
\end{itemize}
\end{tcolorbox}

\vspace{5mm}

%------------------------------------------------------------
% PARTIE 1 : QUELQUES SYMBOLES
%------------------------------------------------------------

\textbf{\large 1. Quelques symboles normalisés}

\vspace{3mm}

Complète le tableau suivant en indiquant la signification et le rôle de chaque symbole :

\vspace{3mm}

\renewcommand{\arraystretch}{1}
\begin{tabular}{|>{\centering\arraybackslash}m{4.5cm}|>{\centering\arraybackslash}m{5cm}|>{\centering\arraybackslash}m{5.5cm}|}
\hline
\rowcolor{blue!15}
\textbf{Symbole} & \textbf{Signification} & \textbf{Rôle} \\
\hline
% Fil conducteur
\begin{circuitikz}[scale=0.9]
  \draw (0,0) -- (2.5,0);
\end{circuitikz}
& \rule[-1.1cm]{0pt}{1.2cm} & \\
\hline
% Lampe
\begin{circuitikz}[scale=0.9]
  \draw (0,0) to[lamp] (2.5,0);
\end{circuitikz}
& \rule[-1.1cm]{0pt}{1.2cm} & \\
\hline
% Interrupteur ouvert
\begin{circuitikz}[scale=0.9]
  \draw (0,0) to[closing switch] (2.5,0);
\end{circuitikz}
\newline {\small (interrupteur ouvert)}
& \rule[-1.1cm]{0pt}{1.2cm} & \\
\hline
% Pile
\begin{circuitikz}[scale=0.9]
  \draw (0,0) to[battery1] (2.5,0);
  \node at (0.55,-0.35) {\small $-$};
  \node at (1.95,-0.35) {\small $+$};
\end{circuitikz}
& \rule[-1.1cm]{0pt}{1.2cm} & \\
\hline
% Générateur (plusieurs piles)
\begin{circuitikz}[scale=1.5]
  \draw (0,0) -- (0.30,0);
  \draw (0.65,0) circle (0.35);
  \node at (0.65,0) {\small\textbf{G}};
  \draw (1.0,0) -- (1.30,0);
\end{circuitikz}
& \rule[-1.1cm]{0pt}{1.2cm} & \\
\hline
% Résistance
\begin{circuitikz}[scale=0.9]
  \draw (0,0) to[R] (2.5,0);
\end{circuitikz}
& \rule[-1.1cm]{0pt}{1.2cm} & \\
\hline
% Moteur
\begin{circuitikz}[scale=1.5]
  \draw (0,0) -- (0.30,0);
  \draw (0.65,0) circle (0.35);
  \node at (0.65,0) {\small\textbf{M}};
  \draw (1.0,0) -- (1.30,0);
\end{circuitikz}
& \rule[-1.1cm]{0pt}{1.2cm} & \\
\hline
\end{tabular}

\vspace{5mm}

\textbf{Appareils de mesure :}

\vspace{3mm}

\renewcommand{\arraystretch}{1}
\begin{tabular}{|>{\centering\arraybackslash}m{4.5cm}|>{\centering\arraybackslash}m{5cm}|>{\centering\arraybackslash}m{5.5cm}|}
\hline
\rowcolor{blue!15}
\textbf{Symbole} & \textbf{Signification} & \textbf{Rôle (ce qu'il mesure)} \\
\hline
\begin{circuitikz}[scale=0.9]
  \draw (0,0) to[ammeter] (2.5,0);
\end{circuitikz}
& \rule[-1.1cm]{0pt}{1.2cm} & \\
\hline
\begin{circuitikz}[scale=0.9]
  \draw (0,0) to[voltmeter] (2.5,0);
\end{circuitikz}
& \rule[-1.1cm]{0pt}{1.2cm} & \\
\hline
\end{tabular}

\newpage

%------------------------------------------------------------
% PARTIE 2 : SCHÉMATISATION
%------------------------------------------------------------

\textbf{\large 2. Schématiser des circuits électriques}

\vspace{4mm}

\noindent
\begin{tabular}{|p{0.47\linewidth}|p{0.47\linewidth}|}
\hline
\textbf{Circuit A} \hfill \textit{(2 pts)} &
\textbf{Circuit B} \hfill \textit{(2 pts)} \\[2mm]
\small Schématiser un circuit en \textbf{série} comprenant : un générateur, deux lampes, un ampèremètre, et un voltmètre branché en dérivation sur l'une des deux lampes. &
\small Schématiser un circuit en \textbf{dérivation} comprenant : un générateur, une lampe, un moteur, et un voltmètre branché en dérivation sur le moteur. \\[2mm]
\hline
\rule{0pt}{7.5cm} & \\
\hline
\end{tabular}

\vspace{6mm}

\noindent
\begin{tabular}{|p{0.47\linewidth}|p{0.47\linewidth}|}
\hline
\textbf{Circuit C} \hfill \textit{(2 pts)} &
\textbf{Circuit D} \hfill \textit{(2 pts)} \\[2mm]
\small Schématiser un circuit en \textbf{série} comprenant : un générateur, deux lampes, et un interrupteur commandant les deux lampes. &
\small Schématiser un circuit en \textbf{dérivation} comprenant : un générateur, deux lampes, et un interrupteur commandant les deux lampes. \\[2mm]
\hline
\rule{0pt}{7.5cm} & \\
\hline
\end{tabular}

\newpage

%------------------------------------------------------------
% PARTIE 3 : RÉALISATION ET OBSERVATIONS
%------------------------------------------------------------

\textbf{\large 3. Réaliser les montages et observer}

\vspace{4mm}

\noindent
\begin{tabular}{|p{0.47\linewidth}|p{0.47\linewidth}|}
\hline
\textbf{Réaliser le circuit C (série)} \hfill \textit{(2 pts)} &
\textbf{Réaliser le circuit D (dérivation)} \hfill \textit{(2 pts)} \\
\hline
\end{tabular}

\appelprof

\vspace{5mm}

\textbf{Questions d'observation :} \hfill \textit{(4 pts)}

\vspace{4mm}

\textbf{1.} Comment varie l'éclairement des deux ampoules dans le circuit en \textbf{série} (C) par rapport au circuit en \textbf{dérivation} (D) ?

\cadreblanc{30}

\textbf{2.} Que se passe-t-il si une ampoule est \textbf{défectueuse} (grillée) dans le montage en \textbf{série} ? Et dans le montage en \textbf{dérivation} ? Explique pourquoi.

\cadreblanc{40}

\appelprof

\vspace{5mm}

\begin{center}
  \textbf{\large Total : /16 points}
\end{center}

\end{document}
